\documentclass[10pt,a4paper]{amsart}
\usepackage[margin=19mm]{geometry}
\usepackage{amsmath,amssymb,mathtools}
\usepackage[scr=rsfs,cal=euler]{mathalfa}
\usepackage{xfrac}
\usepackage[unicode,hidelinks,bookmarks=false]{hyperref}
\newcommand{\ie}{i.e.\ }
\newcommand{\cf}{cf.\ }
\newcommand{\resp}{respectively\ }
\newcommand{\Subsection}{\subsection}
\providecommand{\nobreakdash}{\nobreak\hbox{-}}
\setlength{\emergencystretch}{2em}
\multlinegap0pt

\theoremstyle{plain}
\newtheorem{theorem}{Theorem}
\newtheorem{lemma}[theorem]{Lemma}
\newtheorem{remark}[theorem]{Remark}
\theoremstyle{definition}
\newtheorem{definition}{Definition}[section]
\title[On the distribution of the minimal length of addition chains]{On the distribution of the minimal length of addition chains}
\author{Jean-Marie De Koninck, Nicolas Doyon, William Verreault}
\date{Source: arXiv:2609.28374v1 (CC BY 4.0)}
\begin{document}
\maketitle

\noindent\emph{Source and licence.} The delivered work is the article shown in the title above, version arXiv:2609.28374v1 (CC BY 4.0), distributed under the Creative Commons Attribution 4.0 International licence, as recorded in the arXiv licence metadata for this exact version and shown on the abstract page saved with this item. The full licence notice, which carries the licence URL and the address of that abstract page, is the bundled file \texttt{licenses/arxiv-2609.28374-CC-BY-4.0.txt}.
\par\smallskip

\noindent\textit{Editorial scope.} Counted unit: the article's Lemma (source label lem:F-new-marker-separation) (source0.tex lines 1557--1567, source label \texttt{lem:F-new-marker-separation}): ``If are two admissible triples at the beginning of the same phase, then some coordinate satisfies 2 (b s(p,q,e)-b s(p',q',e') ) 4L. Here is the -adic valuation of , with .''. It is delivered together with the article's complete original proof (source0.tex lines 1569--1612, one proof environment adjacent to the statement, 44 lines), copied line for line from the article's own text. Required context retained uncounted: eq:F-new-ancestry-expansion (lines 1502--1509). Editorial formatting only: an AMS \texttt{amsart} preamble that restates the article's own macro definitions and its own theorem declarations, and the theorem environments the retained lines use; the article's own class file is neither shipped nor input; the retained environments are renumbered sequentially in order of appearance, whereas the article prints them with its own numbering; the article's equation numbering is not reproduced and every cross-reference inside the retained text resolves to the wrapper's numbering. No citation occurs inside any retained line, so the wrapper carries no bibliography; All retained bodies are exact copies of the bundled \texttt{sources/211584/source0.tex}, so no omission receipt applies. No asset-inclusion command occurs inside any retained span and the package ships no image asset.


\section*{Retained context: eq:F-new-ancestry-expansion (source0.tex lines 1502--1509)}

\begin{equation}
 \mathbf a_v
 \equiv
 \mathbf e_r+
 \sum_{x\in\mathcal P(v)}2^{e_x}\mathbf e_{s_x}
 \pmod {2^{2L}}.
 \label{eq:F-new-ancestry-expansion}
\end{equation}


\section{The counted unit: Lemma (source label lem:F-new-marker-separation) and its complete original proof}

\begin{lemma}
\label{lem:F-new-marker-separation}
If $(p,q,e)\neq(p',q',e')$ are two admissible triples at the beginning of the
same phase, then some coordinate $s$ satisfies
\begin{equation}
 \nu_2\bigl(b_s(p,q,e)-b_s(p',q',e')\bigr)<4L.
 \label{eq:F-new-triple-separation}
\end{equation}
Here $\nu_2(x)$ is the $2$-adic valuation of $x$, with
$\nu_2(0)=+\infty$.
\end{lemma}

\begin{proof}
We first note that if $v\neq w$ belong to the same reservoir and have the same
root $r$, then some coordinate distinguishes $\mathbf a_v$ and $\mathbf a_w$
at valuation less than $2L$.  This situation does not arise in
$\mathcal H_0$, which has one seed of every root.  In a later reservoir, let
$r\to s$ be the arrow created when $v$ was born.  The elements $v$ and $w$
were produced in the same preceding phase from old sources.  Hence the arrow
$r\to s$ occurs neither in the principal ancestry of $w$ nor as the arrow
created when $w$ was born: the former ancestry predates that phase, and the
latter arrow is different by the rule forbidding reuse of an ordered pair.
At coordinate $s$,
\eqref{eq:F-new-ancestry-expansion} therefore gives
$$
 a_v(s)-a_w(s)\equiv2^{e_v}\pmod {2^{2L}},
$$
where $L\leq e_v<2L$.  Thus
$$
 \nu_2\bigl(a_v(s)-a_w(s)\bigr)=e_v<2L.
$$
Once created, the arrow $r\to s$ remains used at the beginning of every later
phase.

We now compare the two triples.  If their principal roots differ, their
coefficient vectors already differ modulo $2$.  Suppose their principal roots
agree but $p\neq p'$.  By the preceding paragraph, some previously used arrow
$r\to s$ gives a distinguishing coordinate of valuation less than $2L$.
Admissibility forbids both secondary roots from being $s$.  At coordinate $s$,
$\mathbf a_q$ and $\mathbf a_{q'}$ are therefore divisible by $2^L$, and after
multiplication by $2^e$ and $2^{e'}$, the two new contributions are divisible
by $2^{2L}$ and cannot cancel the old distinction.

It remains to consider $p=p'$.  Suppose first that the secondary roots differ,
and let $s=\operatorname{root}(H_q)$.  At coordinate $s$, the contribution
$2^e\mathbf a_q$ has valuation $e$, whereas
$2^{e'}\mathbf a_{q'}$ is divisible by $2^{e'+L}$.  Since
$e'+L\geq2L>e$, this coordinate distinguishes the two vectors at valuation
$e<2L$.  If the secondary roots agree but $e\neq e'$, the difference at their common root coordinate has valuation $\min(e,e')<2L$.  Finally, if the secondary roots and exponents agree
but $q\neq q'$, the preceding same-root observation, followed by multiplication
by $2^e$, gives a distinguishing coordinate of valuation at most
$$
 (2L-1)+(2L-1)=4L-2.
$$
These cases prove \eqref{eq:F-new-triple-separation}.
\end{proof}

\end{document}
